\chapter{Introduction}\label{sec:intro}

Spherical spin glasses are mean-field models of disordered systems, in which the energy is a random Gaussian function on a high-dimensional sphere. Following the work of Kirkpatrick, Thirumalai, and Wolynes \cite{KirkpatrickThirumalai1987,KirkpatrickWolynes1987,KirkpatrickThirumalaiWolynes1989}, $p$-spin models have served as mean-field models of the glass transition, and spherical models are among those that can be solved: Crisanti and Sommers computed their statics \cite{CrisantiSommers1992}, and Crisanti, Horner, and Sommers their dynamics \cite{CrisantiHornerSommers1993}. For pure spherical $p$-spin models with $p\ge3$, the standard picture involves a dynamical transition temperature $T_d$ and a lower static transition temperature $T_s$. Between $T_s$ and $T_d$, the Langevin dynamics is predicted to fall out of equilibrium, although the free energy keeps its annealed value: equilibrium dynamics is predicted to remain for very long times in one of exponentially many states that together carry the Gibbs measure \cite{CrisantiHornerSommers1993,FranzParisi1995,BarratBurioniMezard1996}. Below $T_s$, the Gibbs measure itself condenses onto few states. Started instead from a random configuration, the dynamics below $T_d$ is predicted to age and to stay away from the states that carry the Gibbs measure, on all time scales that survive the limit of infinite dimension \cite{CugliandoloKurchan1993}. Mixtures can behave differently: the two transitions may coincide (\cref{ex:inf-degenerate}), and the long-time behavior can depend on the preparation \cite{FolenaFranzRicciTersenghi2020}.

In this article, we prove versions of these predictions in finite dimension $N$, on time scales exponential in $N$. We study the Langevin dynamics of spherical mixed $p$-spin glasses at a fixed temperature and ask two questions. The first concerns temperatures below $T_s$: does the dynamics started from the uniform law on the sphere, which does not depend on the disorder, reach equilibrium, and on what time scale? The second concerns the interval $T_s<T<T_d$: what does the equilibrium dynamics remember of its starting point? For the first question, we show that the law of the dynamics started from the uniform law stays at total variation distance tending to one from the Gibbs measure up to a time exponential in $N$. This holds at every temperature below $T_s$ for pure models of degree at least four, and for the mixtures in a class defined by an explicit scalar condition (\cref{thm:main-condensed,cor:main-classes}), and it extends to every initial law chosen independently of the disorder (\cref{thm:main-initial}). For the second question, we show for general finite mixtures that the generator has exponentially many exponentially small eigenvalues, that a fixed part of the spin lies in the corresponding slow modes, that the stationary spin correlation stays above a plateau up to an exponential time, that the dynamics started from most points of the Gibbs measure stays far from equilibrium, and that the two endpoints of the stationary dynamics share mutual information at least linear in $N$ (\cref{thm:main-memory,thm:main-plateau,thm:main-cusp}).

\section{The model and the two thresholds}\label{sec:intro-model}

Let $S_N=\{x\in\R^N:|x|^2=N\}$, and let $\nu_N$ be the uniform probability measure on $S_N$. For $x,y\in\R^N$, we write $R(x,y)=x\cdot y/N$ for the overlap and $U(x)=x/\sqrt N$ for the normalized spin. A \emph{finite mixture} is a nonzero polynomial
\begin{equation*}
\xi(r)=\sum_{p=2}^Pc_pr^p,\qquad c_p\ge0.
\end{equation*}
Its minimum degree $d$ is the smallest $p$ with $c_p>0$, and $\xi$ is \emph{pure} if exactly one coefficient is nonzero. The field $H_N$ is the centered Gaussian process on $S_N$ with covariance
\begin{equation*}
\E H_N(x)H_N(y)=N\xi\bigl(R(x,y)\bigr).
\end{equation*}
For $\xi(r)=r^p$, we may take $H_N(x)=N^{-(p-1)/2}\sum_{i_1,\dots,i_p}g_{i_1\cdots i_p}x_{i_1}\cdots x_{i_p}$ with independent standard Gaussian coefficients, and a mixture is a sum of independent fields of this form with weights $\sqrt{c_p}$. We allow a quadratic term and odd degrees. We refer to the randomness of $H_N$ as the disorder. For an inverse temperature $\beta\ge0$, the Gibbs measure and the free energy are
\begin{equation*}
\pi(\dd x)=\frac{e^{\beta H_N(x)}}{Z_{N,\beta}}\,\nu_N(\dd x),\qquad Z_{N,\beta}=\int_{S_N}e^{\beta H_N}\dd\nu_N,\qquad F(\beta)=\lim_{N\to\infty}\frac1N\E\log Z_{N,\beta}.
\end{equation*}
The Gibbs measure favors large values of $H_N$; in the physics literature, the energy is usually $-H_N$. The limit $F(\beta)$ exists by the Parisi formula, which we recall below.

The Langevin dynamics is the diffusion on $S_N$ with generator
\begin{equation*}
L=\Delta_{S_N}+\beta\nabla_{S_N}H_N\cdot\nabla_{S_N},
\end{equation*}
that is, the solution of $\dd X_t=\sqrt2\,P_{X_t}\dd W_t+\beta\nabla_{S_N}H_N(X_t)\dd t-(1-1/N)X_t\dd t$, where $W$ is a standard Brownian motion in $\R^N$ and $P_x$ is the orthogonal projection onto the tangent space of $S_N$ at $x$. It is reversible with respect to $\pi$, and its Dirichlet form is $\cE(f,f)=\int|\nabla_{S_N}f|^2\dd\pi$. The operator $-L$ has discrete spectrum $0=\lambda_0<\lambda_1\le\lambda_2\le\cdots$, where eigenvalues are counted with multiplicity. We write $P_t=e^{tL}$, $P_t(x,\cdot)$ for its transition kernel, and $\mu P_t$ for the law at time $t$ of the diffusion with initial law $\mu$. All these objects depend on the realization of $H_N$, which we suppress from the notation. The total variation distance is $\norm{\mu-\nu}_{\TV}=\sup_A|\mu(A)-\nu(A)|$. Physicists often write the generator as $T\Delta-\nabla E\cdot\nabla$ with $E=-H_N$; this changes the time scale by a constant factor, which does not matter on the exponential time scales considered here.

The two thresholds are defined through the functions
\begin{equation*}
A(r)=\frac{(1-r)\xi'(r)}r,\qquad\Phi_\beta(r)=\beta^2\xi(r)+r+\log(1-r),\qquad0<r<1.
\end{equation*}
The \emph{dynamical threshold} $\bd=\bd(\xi)$ and the \emph{static threshold} $\bs=\bs(\xi)$ are given by
\begin{equation}\label{eq:intro-thresholds}
\bd^{-2}=\sup_{0<r<1}A(r),\qquad\bs^2=\inf_{0<r<1}\frac{-r-\log(1-r)}{\xi(r)}.
\end{equation}
For a pure model, we write $\bd(p)$ and $\bs(p)$ for the thresholds of $\xi(r)=r^p$.

The static threshold is the point where the free energy leaves its annealed value. By Jensen's inequality, we have $F(\beta)\le\lim_NN^{-1}\log\E Z_{N,\beta}=\beta^2\xi(1)/2$, and equality holds if and only if $\Phi_\beta\le0$ on $(0,1)$, that is, if and only if $\beta\le\bs$. For even mixtures, Talagrand proved that the corresponding strict inequality is sufficient for this equality and that the non-strict inequality is necessary \cite[Proposition~2.3]{Talagrand2006}. The equivalence follows from the optimality criterion \cite[Proposition~2.1]{Talagrand2006} and the Parisi formula, which Chen proved for general mixtures \cite{Chen2013} (see \cref{lem:pre-RS}). Following Huang and Sellke \cite[Definition~3]{HuangSellke2024}, we say that $(\xi,\beta)$ is \emph{strictly replica symmetric} if $\Phi_\beta<0$ on $(0,1)$ (\cref{def:pre-thresholds}). For every $\beta<\bs$, the pair is strictly replica symmetric and $\beta^2\xi''(0)<1$ (\cref{lem:pre-RS}). We call the range $\beta>\bs$ the \emph{condensed phase}.

The dynamical threshold enters through the derivative
\begin{equation*}
\Phi_\beta'(r)=\frac r{1-r}\bigl(\beta^2A(r)-1\bigr).
\end{equation*}
If $\beta\le\bd$, then this identity implies that $\Phi_\beta$ is nonincreasing, hence $\Phi_\beta\le0$. We conclude that $\bd\le\bs$. For $\beta>\bd$, the set $\{r\in(0,1):\beta^2A(r)>1\}$ is nonempty and its closure does not contain $1$, since $A(1^-)=0$. We define its right endpoint
\begin{equation}\label{eq:intro-qstar}
q_*=\sup\bigl\{r\in(0,1):\beta^2(1-r)\xi'(r)>r\bigr\}.
\end{equation}
By continuity, $q_*$ solves the \emph{plateau equation}
\begin{equation}\label{eq:intro-plateau}
\beta^2(1-q)\xi'(q)=q.
\end{equation}
It is the largest solution, unless $\beta^2A$ touches the value $1$ at some point to the right of $q_*$ without exceeding it. In the physics literature, for $\bd<\beta<\bs$, the largest solution of \eqref{eq:intro-plateau} is the predicted long-time limit of the equilibrium spin correlation (see \cref{sec:intro-physics}); below the static transition, the equilibrium plateau involves the parameters of the condensed phase instead. The threshold $\bd$ is the infimum of the inverse temperatures at which \eqref{eq:intro-plateau} has a solution in $(0,1)$. For $\xi(r)=r^p$ with $p\ge3$, the function $A$ is maximized at $q_d=(p-2)/(p-1)$, and $\bd(p)^{-2}=p(p-2)^{p-2}/(p-1)^{p-1}$; for $p=2$, the supremum of $A$ is approached as $r\downarrow0$, and $\bd(2)=\bs(2)=1/\sqrt2$. In the normalization $\xi(r)=r^p/2$, which is common in the physics literature, this is the classical value of the dynamical temperature \cite[Section~4.5]{CastellaniCavagna2005}. The interval $\bd<\beta<\bs$ between the dynamical and static thresholds is nonempty for pure models with $p\ge3$, and more generally whenever $c_2=0$ (see \cref{rem:inf-gap}). It can be empty: for $\xi(r)=r^2/2+r^4$, we have $\bd=\bs=1$ (\cref{ex:inf-degenerate}).

To describe the condensed phase, we recall the Crisanti--Sommers formula. For a probability measure $\zeta$ on $[0,1]$, let $\hat x_\zeta(t)=\int(1-\max\{s,t\})\,\zeta(\dd s)$. The Parisi formula expresses $F(\beta)$ as the minimum over $\zeta$ of an explicit functional of $\hat x_\zeta$, and its minimizer is the Parisi measure (\cref{def:pre-CS,thm:pre-parisi}). This formula was derived by Crisanti and Sommers \cite{CrisantiSommers1992} and proved by Talagrand for even mixtures \cite[Theorem~1.1]{Talagrand2006} and by Chen in general \cite[Theorem~1]{Chen2013}; Huang and Sellke gave a constructive proof \cite{HuangSellke2024}. The Parisi measure is $\delta_0$ exactly when $\beta\le\bs$. A minimizer $\zeta$ is characterized by its first variation: the function
\begin{equation*}
V_\zeta(r)=\int_0^r\frac{r-v}{\hat x_\zeta(v)^2}\dd v-\beta^2\xi(r),\qquad0\le r<1,
\end{equation*}
is nonnegative and vanishes on the support of $\zeta$ \cite[Proposition~2.1]{Talagrand2006}, \cite[Lemma~1.2]{HuangSellke2024}. For the Parisi measure $\zeta$, we call $V_\zeta$ the \emph{first-variation function} and its zero set in $[0,1)$ the \emph{contact set}. The contact set contains the support of $\zeta$, and it can be larger \cite[Remark~1.4]{HuangSellke2024}. We consider measures with two atoms, $\zeta=x\delta_0+(1-x)\delta_Q$, which describe one-step replica symmetry breaking (1RSB). For such $\zeta$, $\hat x_\zeta$ is the function $\hat x$ of the next definition, and $V_\zeta$ is the function $V$.

\begin{definition}[Strict one-step replica symmetry breaking]\label{def:intro-certificate}
Let $\xi$ be a finite mixture with minimum degree at least four, and let $\beta>0$. For $x,Q\in(0,1)$, let $s=1-Q$, $C=s+xQ$, and
\begin{equation*}
\hat x(v)=\begin{cases}C-xv,&0\le v\le Q,\\1-v,&Q\le v<1,\end{cases}
\qquad
V(r)=\int_0^r\frac{r-v}{\hat x(v)^2}\dd v-\beta^2\xi(r),\quad0\le r<1.
\end{equation*}
The pair $(\xi,\beta)$ is \emph{strictly 1RSB} if there exist $x,Q\in(0,1)$ such that its Parisi measure is $x\delta_0+(1-x)\delta_Q$ and its contact set $\{r\in[0,1):V(r)=0\}$ is exactly $\{0,Q\}$. The contact at $Q$ is \emph{nondegenerate} if $V''(Q)>0$. Equivalently (\cref{sec:th-certificate}), $(\xi,\beta)$ is strictly 1RSB with a nondegenerate contact at $Q$ if and only if there exist $x,Q\in(0,1)$ such that
\begin{equation}\label{eq:intro-certificate}
V(Q)=V'(Q)=0,\qquad V(r)>0\quad\text{for }r\in(0,1)\setminus\{Q\},\qquad V''(Q)>0.
\end{equation}
\end{definition}

Note that $V(0)=V'(0)=0$ holds automatically. The conditions \eqref{eq:intro-certificate} involve one explicit function of one variable, and they can be checked without the Parisi formula; we call them the \emph{1RSB assumptions}. They imply that $x\delta_0+(1-x)\delta_Q$ is the Parisi measure, so that $x$ and $Q$ are determined by $(\xi,\beta)$ (\cref{prop:th-structure}). Since $\hat x(Q)=1-Q$, the condition $V''(Q)>0$ is equivalent to $\beta^2(1-Q)^2\xi''(Q)<1$, that is, to the positivity of the replicon eigenvalue $(1-Q)^{-2}-\beta^2\xi''(Q)$ at $Q$ \cite[Section~1.3]{BenArousJagannath2018}. In the language of physics, the Gibbs measure is one-step replica symmetry breaking: $Q$ is the overlap of two samples from the same state, and $x$ is the mass of the atom of the Parisi measure at zero. Our terminology follows Huang and Sellke, who call a model strictly replica symmetric if its contact set is $\{0\}$ \cite[Definition~3]{HuangSellke2024} and, at zero temperature, strictly 1RSB if the corresponding contact set is $\{0,1\}$ \cite[Definition~4]{HuangSellke2024}; \cref{def:intro-certificate} is the analogous notion at positive temperature. Every pair that satisfies it lies in the condensed phase, $\beta>\bs$ (\cref{lem:th-superstatic}).

\section{The uniform initial law and Gibbs-typical starting points}\label{sec:intro-starts}

The two questions of this paper concern two different initial laws, and the distinction is essential. Three ways of measuring the distance to equilibrium at time $t$ are
\begin{equation*}
\sup_{x\in S_N}\norm{P_t(x,\cdot)-\pi}_{\TV},\qquad\norm{\nu_NP_t-\pi}_{\TV},\qquad\int\norm{P_t(x,\cdot)-\pi}_{\TV}\,\pi(\dd x).
\end{equation*}
An exponentially small spectral gap $\lambda_1$ means that the first quantity decays slowly: some observable relaxes slowly from some starting points. It does not show that the dynamics started from $\nu_N$ excites this observable, since the second quantity depends on the projection of the density $\dd\nu_N/\dd\pi$ onto the slow modes. Gheissari, Grant, and Yu gave an example of this phenomenon \cite[Theorem~1.2]{GheissariGrantYu2026}.
In a spiked spherical matrix model at low temperature, the worst-case mixing time is exponential in $N$, while the dynamics started from any initial law that is symmetric in the coordinate along the top eigenvector, including the uniform law, mixes in time of order $\log N$. The exponentially small spectral gaps that are known for spherical spin glasses (see \cref{sec:intro-math}) therefore do not answer the first question.

The uniform law $\nu_N$ is the Gibbs measure at infinite temperature, and it does not depend on the disorder. The dynamics started from $\nu_N$ describes a quench from infinite temperature to the temperature $1/\beta$. This is the protocol of Cugliandolo and Kurchan \cite{CugliandoloKurchan1993}, whose initial configuration is chosen independently of the disorder. For brevity, we call $\nu_N$ the \emph{uniform start}.

If instead the initial point is drawn from $\pi$, then $\pi P_t=\pi$ for every $t$, and the law of $X_t$ is in equilibrium at all times. What the equilibrium dynamics can remember lies in the joint law of $X_0$ and $X_t$. It is visible in the conditional law $P_t(X_0,\cdot)$ of the dynamics started from a point $X_0$ drawn from $\pi$ and then held fixed, in the stationary correlation $\E_\pi R(X_0,X_t)$, and in the mutual information between $X_0$ and $X_t$. We speak of \emph{Gibbs-typical starting points}. They depend on the disorder, and the example $\mu=\pi$ shows that no separation result for the uniform start can hold for all initial laws that depend on the disorder. \Cref{thm:main-initial} below says how much dependence is allowed.

\section{Main results}\label{sec:intro-results}

Throughout, the covariance $\xi$ and the inverse temperature $\beta$ are fixed before $N\to\infty$, and probabilities refer to the disorder. The conventions for constants and for statements that hold with high probability are fixed in \cref{sec:pre-conventions}.

\subsection{The condensed phase}\label{sec:intro-condensed}

Our first result concerns the dynamics started from the uniform law in the condensed phase.

\begin{theorem}[Slow equilibration from the uniform law]\label{thm:main-condensed}
Let $\xi$ be a finite mixture with minimum degree $d\ge4$, and let $\beta>0$ be such that $(\xi,\beta)$ is strictly 1RSB with a nondegenerate contact at $Q$. Then there exists $\gamma=\gamma(\xi,\beta)>0$ such that
\begin{equation}\label{eq:intro-main}
\inf_{0\le t\le e^{\gamma N}}\norm{\nu_NP_t-\pi}_{\TV}\longrightarrow1\qquad\text{in probability as }N\to\infty.
\end{equation}
\end{theorem}

We make no assumption on the number of degrees with $c_p>0$, on their parity, or on the sizes of the coefficients. By the contraction of total variation along a Markov semigroup (\cref{lem:pre-contraction}), the infimum in \eqref{eq:intro-main} is the distance at the final time $e^{\gamma N}$. In fact, the proof shows that one set works for all times. For every $\eps>0$ and all sufficiently large $N$, with probability at least $1-\eps$, there exists a set $B$, depending on the disorder, such that $\pi(B)\ge1-\eps$ and $(\nu_NP_t)(B)\le\eps$ for all $t\le e^{\gamma N}$. One may take for $B$ the set $B_{N,M}$ of \cref{def:red-centers}: the union of the bands of a small fixed width around the critical points $u$ of the center field $H_Q(u)=H_N(\sqrt Q\,u)$ with $H_Q(u)\in[m_N-M,m_N+M]$ and with normalized radial derivative $K_Q(u)/N$ in a fixed interval, where $K_Q(u)$ is the derivative of $t\mapsto H_N(t\sqrt Q\,u)$ at $t=1$, $m_N$ is a deterministic centering, and $M$ depends on $\eps$. With probability close to one, these values lie within a bounded distance of the maximum of $H_Q$ (see \cref{sec:intro-ideas}).

The 1RSB assumptions can be checked in many cases. The following corollary lists the classes for which we verify them.

\begin{corollary}[Strictly 1RSB models]\label{cor:main-classes}
The conclusion of \cref{thm:main-condensed} holds in each of the following cases.
\begin{enumerate}[label=(\alph*)]
\item The pure model $\xi(r)=r^p$, for every $p\ge4$ and every $\beta>\bs(p)$.
\item Every finite mixture $\xi$ with minimum degree at least four such that $r\mapsto\xi''(r)^{-1/2}$ is strictly convex on $(0,1]$, and every $\beta>\bs(\xi)$. This holds in particular if the smallest and the largest degrees $d$ and $P$ with $c_p>0$ satisfy $6(d-2)(P-2)>(d+P-5)^2$; for example, if all degrees with $c_p>0$ lie in $\{4,5,\dots,11\}$.
\item The covariances $\xi(r)=r^d+\eps\sum_{d<p\le P,\ p\text{ even}}b_pr^p$, where $d\ge4$ is even, $P>d$, $\beta>\bs(d)$, the numbers $b_p\ge0$ satisfy $\sum_pb_p=1$, and $0<\eps<\eps_0$, for some $\eps_0=\eps_0(d,P,\beta)>0$ that does not depend on $(b_p)$.
\item The covariance $\xi(r)=r^4+br^{12}$ at the inverse temperature $\beta$, where
\begin{equation*}
t=\frac{45-24\log6}{72\log6-85},\qquad b=t\Bigl(\frac{10}9\Bigr)^8,\qquad\beta^2=\frac{15000}{729(4+12t)}.
\end{equation*}
\end{enumerate}
\end{corollary}

Numerically, $b\approx0.10546$ and $\beta\approx2.12778$ in (d), and the 1RSB assumptions hold there with $Q=9/10$ and $x=5/9$. In case (d), the function $\xi''^{-1/2}$ is not convex on $(0,1]$, so (d) is not contained in (b) (\cref{prop:th-example}). In case (b), the degrees may have either parity, as in (a). The pure cubic model and models with a quadratic or cubic component are not covered; we return to this point in \cref{sec:intro-scope}.

Next, we transfer \cref{thm:main-condensed} to other initial laws, with the same exponent $\gamma$.

\begin{theorem}[Other initial laws]\label{thm:main-initial}
Let $\xi$ and $\beta$ satisfy the hypotheses of \cref{thm:main-condensed}, let $\gamma$ be the exponent given there, and let $T_N=e^{\gamma N}$.
\begin{enumerate}[label=(\roman*)]
\item For every $N$, every Borel probability measure $\mu$ on $S_N$, and every $T\ge0$,
\begin{equation}\label{eq:intro-haar}
\E\bigl[1-\norm{\mu P_T-\pi}_{\TV}\bigr]\le\E\bigl[1-\norm{\nu_NP_T-\pi}_{\TV}\bigr].
\end{equation}
In particular, for every $\eps>0$,
\begin{equation*}
\sup_\mu\Pp\Bigl(\inf_{0\le t\le T_N}\norm{\mu P_t-\pi}_{\TV}<1-\eps\Bigr)\longrightarrow0,
\end{equation*}
where the supremum is over all Borel probability measures $\mu$ on $S_N$. The same convergence holds for every random initial law that is independent of $H_N$, in the joint probability of the disorder and the random initial law.
\item Let $0\le\kappa<(d-2)/2$ and $C_0>0$. For each $N$, let $\mu_{N,1},\dots,\mu_{N,K_N}$ be Borel probability measures on $S_N$ that do not depend on $H_N$, with $K_N\le C_0N^\kappa$. Then
\begin{equation*}
\inf\Bigl\{\Bigl\|\Bigl(\sum_{j=1}^{K_N}w_j\mu_{N,j}\Bigr)P_t-\pi\Bigr\|_{\TV}:\ w_j\ge0,\ \sum_{j=1}^{K_N}w_j=1,\ 0\le t\le T_N\Bigr\}\longrightarrow1
\end{equation*}
in probability. If the collection is random and independent of $H_N$, the same convergence holds in the joint probability.
\item Let $H\mapsto\mu^H_N$ be a Borel probability kernel from fields to $S_N$ such that $\KL(\mu^{H_N}_N\Vert\nu_N)/\log N\to0$ in probability. Then
\begin{equation*}
\inf_{0\le t\le T_N}\norm{\mu^{H_N}_NP_t-\pi}_{\TV}\longrightarrow1\qquad\text{in probability.}
\end{equation*}
\end{enumerate}
\end{theorem}

In part (i), we may take any deterministic sequence of initial laws, for instance point masses $\mu_N=\delta_{x_N}$, and the bound on the probability of failure is uniform in the initial law. The supremum is outside the probability, and it cannot be moved inside: the initial law $\mu=\pi$ depends on the disorder and is at distance zero from $\pi$ at all times. In part (ii), the infimum over the weights allows them to depend on the disorder. For example, among $N^\kappa$ prescribed points, one may choose, after observing the disorder, the point from which the dynamics comes closest to equilibrium. Such a choice can carry information of order $\kappa\log N$ about the disorder. In part (iii), the initial law may depend on the disorder in an arbitrary measurable way, provided that its relative entropy with respect to $\nu_N$ is $o(\log N)$; for comparison, the relative entropy of $\pi$ with respect to $\nu_N$ is of order $N$.

\subsection{Below the static threshold}\label{sec:intro-prestatic}

We now turn to the interval $\bd<\beta<\bs$ and to the equilibrium dynamics. Recall the plateau value $q_*$ from \eqref{eq:intro-qstar}. The centered spin $Z_N(x)=U(x)-\pi(U)$ is an $\R^N$-valued function on $S_N$, and operators act on it componentwise.

\begin{theorem}[Slow relaxation of the equilibrium dynamics]\label{thm:main-memory}
Let $\xi$ be a finite mixture with $\bd(\xi)<\bs(\xi)$, let $\beta\in(\bd,\bs)$, and let $q_*=\sup\{r\in(0,1):\beta^2(1-r)\xi'(r)>r\}$. For every $0<b<q_*^2$, there exist $\kappa,\gamma,a,c,C>0$, depending on $\xi$, $\beta$, and $b$, such that for all sufficiently large $N$, with probability at least $1-Ce^{-aN}$, the following hold.
\begin{enumerate}[label=(\roman*)]
\item $\lambda_{\lfloor e^{\kappa N}\rfloor}\le e^{-2\gamma N}$.
\item With $Z_N=U-\pi(U)$ and $\Pi_N=\ind_{(0,e^{-2\gamma N}]}(-L)$, acting componentwise, we have $\norm{\Pi_NZ_N}_{L^2(\pi;\R^N)}^2\ge b$.
\item For the stationary diffusion, with $X_0\sim\pi$, we have $\E_\pi R(X_0,X_t)-|\pi(U)|^2\ge b$ for all $0\le t\le e^{\gamma N}$.
\item There exists a Borel set $\cX_N\subset S_N$ with $\pi(\cX_N)\ge1-Ce^{-aN}$ such that $\norm{P_t(x,\cdot)-\pi}_{\TV}\ge1-Ce^{-aN}$ for all $x\in\cX_N$ and $0\le t\le e^{\gamma N}$.
\item For the stationary diffusion in the given field, the mutual information satisfies $I(X_0;X_t)\ge cN$ for all $0<t\le e^{\gamma N}$.
\end{enumerate}
\end{theorem}

We allow a quadratic term and odd degrees, and we make no assumption on the shape of $A$: the endpoint $q_*$ may be a degenerate root of \eqref{eq:intro-plateau}, and $A$ may have several local maxima. All five conclusions hold on one event. In words, the law of the stationary dynamics is $\pi$ at every time, but its endpoints stay strongly dependent for an exponential time. We comment on the conclusions in turn.

Conclusion (i) asserts that exponentially many eigenvalues are exponentially small, and (ii) shows that a fixed part of the centered spin lies in the span of the corresponding eigenfunctions, which we call slow modes. Since $\pi$ is invariant, we have $\E_\pi R(X_0,X_t)-|\pi(U)|^2=\ip{Z_N}{P_tZ_N}_{L^2(\pi;\R^N)}$, and by the spectral theorem this quantity is nonnegative and nonincreasing in $t$. Conclusion (iii) states that this connected correlation stays above $b$ up to time $e^{\gamma N}$; up to a factor $e^{-e^{-\gamma N}}$, it follows from (ii). Conclusion (iv) concerns the dynamics started from Gibbs-typical points: for most starting points $x$ under $\pi$, the law at every time up to $e^{\gamma N}$ is at total variation distance close to one from $\pi$. It concerns starting points that depend on the disorder, and it says nothing about the uniform start. Conclusion (v) states that the endpoint $X_t$ of the stationary dynamics carries information at least $cN$ about the starting point; by the chain identity of \cref{lem:pre-MI}, $I(X_0;X_t)=\int\KL(P_t(x,\cdot)\Vert\pi)\,\pi(\dd x)$.

The plateau level in (ii) and (iii) is $q_*^2$, while the physics literature predicts that the stationary correlation stays close to the largest solution of \eqref{eq:intro-plateau}, which is $q_*$ except at the exceptional temperatures noted after \eqref{eq:intro-plateau}. In the next theorem, we reach every level below $q_*$, at the cost of a weaker probability bound for (ii) and (iii).

\begin{theorem}[Correlation bounds at every level below $q_*$]\label{thm:main-plateau}
Let $\xi$, $\beta$, and $q_*$ be as in \cref{thm:main-memory}. For every $0<b<q_*$, there exist $\kappa,\gamma,a,c,C>0$, depending on $\xi$, $\beta$, and $b$, such that conclusions (ii) and (iii) of \cref{thm:main-memory} hold with probability tending to one, and conclusions (i), (iv), and (v) hold for all sufficiently large $N$ with probability at least $1-Ce^{-aN}$.
\end{theorem}

We do not prove a matching upper bound: the theorem does not show that the correlation stays close to $q_*$, only that it does not fall below any fixed level smaller than $q_*$.

Our last main result describes the slow modes as $\beta$ decreases to $\bd$. It holds for infinite mixtures $\xi(r)=\sum_{p\ge2}c_pr^p$ with $c_p\ge0$ that satisfy
\begin{equation}\label{eq:intro-summability}
\sum_{p\ge2}p^{5/2}\sqrt{c_p}<\infty.
\end{equation}
For such $\xi$, the field $H_N$ is twice continuously differentiable, and the dynamics and its spectrum are defined as before (\cref{lem:inf-field,lem:inf-dynamics}). The thresholds and $A$ are defined by the same formulas, and $A$ extends continuously to $[0,1]$ with $A(0)=\xi''(0)$ and $A(1)=0$.

\begin{theorem}[Small eigenvalues near the dynamical threshold]\label{thm:main-cusp}
Let $\xi$ satisfy \eqref{eq:intro-summability} and $\bd(\xi)<\bs(\xi)$, and let $A_{\max}=\sup_{(0,1)}A=\bd^{-2}$. Then the set of global maximizers of $A$ on $[0,1]$ is a finite subset of $(0,1)$. Let $q_d$ be its largest element. There exist an integer $m\ge1$ and $a>0$ such that
\begin{equation*}
A(q_d+y)=A_{\max}\bigl[1-ay^{2m}+O(y^{2m+1})\bigr]\qquad(y\to0).
\end{equation*}
For $\beta>\bd$, let $\eps=\beta^2A_{\max}-1$, and set
\begin{equation*}
C_d=\frac{2m\,q_d}{(2m+1)(1-q_d)}\,a^{-1/(2m)}.
\end{equation*}
There exists $\kappa_*>0$ such that for every $\vartheta\in(0,1)$, there exists $\eps_0(\vartheta)>0$ such that for every fixed $\eps\in(0,\eps_0(\vartheta))$, there exist $a',C>0$ such that, for all sufficiently large $N$, with probability at least $1-Ce^{-a'N}$,
\begin{equation*}
\lambda_{\lfloor e^{\kappa_*N}\rfloor}\le\exp\bigl[-(1-\vartheta)C_d\,\eps^{1+1/(2m)}N\bigr].
\end{equation*}
\end{theorem}

The number of slow modes stays exponential in $N$ with an exponent $\kappa_*$ that does not depend on $\eps$, while the exponent of the eigenvalue bound tends to zero with $\eps$. The maximizer $q_d$ may be degenerate, and it need not be unique; the order $2m$ of its contact with the maximum is determined by $\xi$. When the maximum is nondegenerate, $m=1$ and the exponent is $3/2$. The quantity $C_d\eps^{1+1/(2m)}$ is the leading order of the height of a barrier in the explicit trial pressure of \cref{lem:pr-trial}, which bounds from below the corresponding barrier of the constrained pressure of \cref{def:pr-FP}, the potential of Franz and Parisi (see \cref{sec:intro-physics}). For the pure cubic model, $q_d=1/2$, $A_{\max}=3/4$, $m=1$, and $C_d=1/3$ (\cref{ex:inf-cubic}). The constants $a'$ and $C$ and the threshold for $N$ may depend on $\eps$ and $\vartheta$.

\subsection{Scope}\label{sec:intro-scope}

We record here what the theorems do not assert. The convergence in \cref{thm:main-condensed} is in probability; we do not show that the probability of failure is exponentially small. The covariance and the temperature are fixed before $N\to\infty$. We do not assert that the exponent $\gamma$ of \cref{thm:main-condensed} is uniform near $\bs$; the uniformity in the temperature is that stated in \cref{thm:main-cusp}. The proof of \cref{thm:main-condensed} requires minimum degree at least four: for a component of degree two or three, an error term in our hitting estimate for the caps around critical points does not tend to zero (\cref{rem:cp-scope}). In particular, the pure cubic model is not covered in the condensed phase, although it is covered by \cref{thm:main-memory,thm:main-plateau,thm:main-cusp}. In \cref{thm:main-initial}(ii), the endpoint $\kappa=(d-2)/2$ is not covered. We do not study the dynamics started from the uniform law for $\bd<\beta\le\bs$, where, for pure spherical models, it is also predicted to be slow \cite[arXiv v1, Section~2.5, Conjecture~6]{BenArousJagannath2024}; see \cref{sec:intro-comparison}. In \cref{thm:main-plateau}, the probability in (ii) and (iii) tends to one without a rate. \Cref{thm:main-cusp} bounds eigenvalues only. Finally, the exponent $\kappa$ in \cref{thm:main-memory}(i) is a lower bound for the growth rate of the number of slow modes; we do not identify it with the complexity of metastable states.

\section{Background from physics}\label{sec:intro-physics}

The distinction between a temperature of dynamical arrest and a lower static, or Kauzmann, temperature goes back to Kirkpatrick and Thirumalai, who connected the dynamical equations of $p$-spin models with theories of the structural glass transition \cite{KirkpatrickThirumalai1987}, and to Kirkpatrick and Wolynes, who explained freezing above the thermodynamic transition through metastable states and their configurational entropy \cite{KirkpatrickWolynes1987}. Kirkpatrick, Thirumalai, and Wolynes developed the corresponding activated picture for liquids in finite dimension \cite{KirkpatrickThirumalaiWolynes1989}. For the spherical $p$-spin model, Crisanti and Sommers computed the statics and found a one-step replica symmetry breaking transition \cite{CrisantiSommers1992}, and Crisanti, Horner, and Sommers analyzed the dynamics and found a dynamical transition above the static one \cite{CrisantiHornerSommers1993}. In our notation, these transitions occur at $\bd$ and $\bs$. A pedagogical account of statics, dynamics, and complexity is given by Castellani and Cavagna \cite{CastellaniCavagna2005}. Between the two transitions, the free energy has its annealed value, but the dynamics is not expected to decorrelate: the replica symmetry of the Gibbs measure does not imply fast relaxation of the equilibrium dynamics.

Between the two transitions, the correlation of the equilibrium dynamics, computed in the limit $N\to\infty$, tends to a nonzero plateau at large times. In our normalization, the plateau solves \eqref{eq:intro-plateau} \cite{CrisantiHornerSommers1993}. For the pure model, this equation appears in the work of Barrat, Burioni, and M\'ezard, who select its largest solution \cite[Equation~(14)]{BarratBurioniMezard1996}, and for mixtures in the work of Folena, Franz, and Ricci-Tersenghi \cite[Equation~(A2)]{FolenaFranzRicciTersenghi2020}.
Below $T_s$, the equilibrium plateau involves the parameters of the condensed phase instead; in our notation, $Q$ solves $\beta^2\xi'(Q)=Q/((1-Q)(1-Q+xQ))$ (\cref{prop:th-structure}). The literature selects the largest solution of \eqref{eq:intro-plateau}, the nonergodicity parameter of the stationary dynamics, which is often called the dynamical Edwards--Anderson parameter. Our $q_*$ agrees with this solution unless $\beta^2A$ touches the value $1$ to the right of $q_*$; at such exceptional temperatures, we use $q_*$ as the reference value for our lower bounds. The stationary plateau should be distinguished from the plateau of the aging dynamics started from a random configuration, which satisfies the marginality condition $\beta^2(1-q)^2\xi''(q)=1$ \cite[Equation~(8)]{FranzParisi1995}. When $q_d$ is an interior maximizer of $A$, the two conditions agree at $\beta=\bd$ and $q=q_d$, but they differ for $\beta>\bd$.

Cugliandolo and Kurchan solved the dynamics started from a configuration chosen independently of the disorder, in the limit where $N\to\infty$ before the time tends to infinity \cite{CugliandoloKurchan1993}. Below $T_d$, this dynamics ages: correlations depend on both times, time translation invariance is lost, and the energy approaches the energy of the marginal threshold TAP states instead of the equilibrium energy. The associated threshold of the Hamiltonian at zero temperature is $2\sqrt{(p-1)/p}$ per site for $\xi(r)=r^p$ in our normalization; the asymptotic energy at positive temperature depends on the temperature. Barrat, Burioni, and M\'ezard contrasted this behavior with the dynamics started from equilibrium, which remains in the state in which it starts \cite{BarratBurioniMezard1996}. Kurchan and Laloux explained the slowness of aging through the geometry of phase space, without trapping in stable metastable states \cite{KurchanLaloux1996}. For mixed models under gradient descent at zero temperature after a preparation at positive temperature, with extensions to small thermal noise, Folena, Franz, and Ricci-Tersenghi found that the asymptotic energy depends on the preparation and that the dynamics keeps a memory of its initial condition, in contrast with the picture of a single attracting threshold in pure models \cite{FolenaFranzRicciTersenghi2020}. Folena and Zamponi subsequently gave evidence against weak ergodicity breaking in classes of mixed models \cite{FolenaZamponi2023}. Citro and Ricci-Tersenghi obtained numerical evidence for persistent memory during aging in the spherical $3+4$ model \cite{CitroRicciTersenghi2025}, and Lang, Sachdev, and Diehl studied a transition between strong and weak ergodicity breaking as the preparation temperature varies \cite{LangSachdevDiehl2025}. These works concern the dynamical mean-field equations and numerical evidence, rather than finite-$N$ probability bounds up to exponential times. Lang, Citro, Leuzzi, and Ricci-Tersenghi studied mixed spherical $(p+s)$-spin models quenched from uniformly random configurations to zero temperature. They solved the asymptotic equations of Cugliandolo and Kurchan under the assumptions of ultrametricity between time scales and local equilibration on each time scale, found aging states with one, two, or continuously many effective temperatures, and obtained the asymptotic energies analytically \cite{LangCitroLeuzziRicciTersenghi2026}.

Franz and Parisi introduced the free energy of configurations constrained to have a prescribed overlap with a reference configuration drawn from equilibrium \cite{FranzParisi1995}; see also \cite{FranzParisi1997}. In the replica-symmetric phase, their potential is $(F(\beta)-F_{\mathrm{FP}}(r))/\beta$, where $F_{\mathrm{FP}}$ is the constrained pressure of \cref{def:pr-FP}, so that a local maximum of $F_{\mathrm{FP}}$ at a positive overlap corresponds to metastable states. Monasson counted metastable states by weakly coupling replicas \cite{Monasson1995}, and Kurchan, Parisi, and Virasoro interpreted saddle points of the replica free energy as metastable states and barriers \cite{KurchanParisiVirasoro1993}. Cavagna, Giardina, and Parisi studied the structure of states, the barriers between them, and the stationary points of the TAP free energy \cite{CavagnaGiardinaParisi1997a,CavagnaGiardinaParisi1997b,CavagnaGiardinaParisi1998}. Barrat, Burioni, and M\'ezard also followed an equilibrium state through a change of temperature \cite{BarratBurioniMezard1996}. Barrat, Franz, and Parisi developed this following of states, including bifurcations \cite{BarratFranzParisi1997}, and Sun, Crisanti, Krzakala, Leuzzi, and Zdeborov\'a revisited it for spherical $s+p$ models with planted configurations \cite{SunCrisantiKrzakalaLeuzziZdeborova2012}. Quiet planting, introduced by Krzakala and Zdeborov\'a for random constraint satisfaction problems \cite{KrzakalaZdeborova2009}, represents a typical equilibrium configuration through a planted law of the disorder; we use a planted law in the same way in \cref{sec:pressure}.

The clustering and condensation transitions were formulated for random constraint satisfaction problems by Krzakala, Montanari, Ricci-Tersenghi, Semerjian, and Zdeborov\'a \cite{KrzakalaMontanariRicciTersenghiSemerjianZdeborova2007}. M\'ezard and Montanari related the dynamical transition to reconstruction on trees \cite{MezardMontanari2006}, and Montanari and Semerjian proved inequalities between point-to-set correlation lengths and relaxation times, with exponential lower bounds on relaxation times for sparse models \cite{MontanariSemerjian2006}. Biroli and Kurchan defined metastable states and their configurational entropy through the low-lying spectrum of the Fokker--Planck operator \cite{BiroliKurchan2001}, and T\u{a}nase-Nicola and Kurchan made this construction explicit \cite{TanaseNicolaKurchan2004}; related ideas appear in the work of Gaveau and Schulman \cite{GaveauSchulman1998}. Biroli derived dynamical TAP equations \cite{Biroli1999}. Lopatin and Ioffe used dynamical instantons to study barriers between metastable states in the spherical $p$-spin model \cite{LopatinIoffe2000}. Ros, Biroli, and Cammarota computed the complexity of barriers near a reference minimum \cite{RosBiroliCammarota2019}, and combined Kac--Rice analysis with dynamical mean-field theory to construct activated trajectories through index-one saddles \cite{RosBiroliCammarota2021}. Rizzo computed the probabilities of activated transitions between states by path integrals \cite{Rizzo2021}. Charbonneau, Folena, Malatesta, Rizzo, and Zamponi studied rare trajectories from equilibrium through a dynamical potential at prescribed final overlap, emphasizing that crossing a typical saddle need not decorrelate the dynamics \cite{CharbonneauFolenaMalatestaRizzoZamponi2026}. Near $T_d$, Franz described the constrained potential by the cubic normal form of a spinodal \cite[Section~5.2.1]{Franz2005}, following the classical theory of nucleation near a spinodal \cite{UngerKlein1984}. In this scalar normal form, the height of the barrier between the local minimum and the saddle point scales as $|T-T_d|^{3/2}$. This is the height of the scalar potential barrier; the spatial droplet barrier studied in \cite[Section~5.2.1]{Franz2005} also includes a dimension-dependent factor.

We now compare our results with these predictions. \Cref{thm:main-condensed} is consistent with the aging picture of \cite{CugliandoloKurchan1993}, in which the dynamics from a disorder-independent start does not reach the states that carry the Gibbs measure. It is a statement of a different kind: it holds in finite dimension up to times $e^{\gamma N}$, but only below the static transition, and it says nothing about the energy reached by the dynamics. Conclusions (iii) of \cref{thm:main-memory,thm:main-plateau} are lower bounds for the stationary correlation at levels below $q_*$, which is the plateau predicted in \cite{CrisantiHornerSommers1993,BarratBurioniMezard1996,FolenaFranzRicciTersenghi2020} except at the exceptional temperatures noted above; we do not prove that the correlation is close to $q_*$. Conclusion (iv) is a finite-dimensional form of the prediction that a configuration drawn from equilibrium remains in its state \cite{BarratBurioniMezard1996,FranzParisi1995}. Conclusion (i) is close to the spectral description of metastable states in \cite{BiroliKurchan2001,TanaseNicolaKurchan2004}, but we do not identify the exponent $\kappa$ with the configurational entropy. Conclusion (v) is an information-theoretic expression of the persistence of the state of the dynamics, related to the configurational entropy \cite{Monasson1995} and to the interpretation of reconstruction in \cite{MezardMontanari2006}. We did not locate an earlier explicit statement of a lower bound, linear in $N$, on the stationary endpoint information up to exponential times in our regime; in the regime of \cite{ElAlaouiMontanariSellke2025}, such a bound follows from the persistence of clusters (see \cref{sec:intro-comparison}). For $m=1$, the exponent $3/2$ in \cref{thm:main-cusp} matches the barrier exponent of the cubic normal form \cite{Franz2005}, and $1+1/(2m)$ is the same computation for a contact of order $2m$. The number of slow modes stays exponential as $\beta\downarrow\bd$, which is consistent with the picture in which the complexity of states need not vanish at $T_d$.

\section{Mathematical background}\label{sec:intro-math}

\subsection{Statics and the geometry of the Gibbs measure}\label{sec:intro-statics}

The free energy of spherical models was computed by Crisanti and Sommers \cite{CrisantiSommers1992}. Talagrand proved their formula for even mixtures \cite{Talagrand2006}, Chen extended it to general mixtures \cite{Chen2013}, and Huang and Sellke gave a constructive proof of the lower bound \cite{HuangSellke2024}. The ground state energy was computed by Chen and Sen \cite{ChenSen2017} and by Jagannath and Tobasco \cite{JagannathTobasco2017low}, and properties of Parisi measures and contact sets of spherical models were studied in \cite{AuffingerChen2015,JagannathTobasco2018contacts,AuffingerChen2018}. The TAP approach was introduced by Thouless, Anderson, and Palmer \cite{ThoulessAndersonPalmer1977}, and it was applied to spherical pure models by Kurchan, Parisi, and Virasoro \cite{KurchanParisiVirasoro1993} and by Crisanti and Sommers \cite{CrisantiSommers1995TAP}. Subag developed a generalized TAP approach for spherical models \cite{Subag2024}. He introduced multisamplable overlaps, values $q$ such that many independent samples from the Gibbs measure can have pairwise overlaps close to $q$, and characterized them by equality in a TAP representation of the free energy \cite[Theorem~5]{Subag2024}. For Ising spins, the generalized TAP free energy was developed by Chen, Panchenko, and Subag \cite{ChenPanchenkoSubag2023,ChenPanchenkoSubag2021}; see also \cite{ChenPanchenko2018}. Subag computed the free energy of pure models at all temperatures rigorously from the generalized TAP approach, without the Parisi formula \cite{Subag2023pure}; the resulting formula had been predicted in \cite{KurchanParisiVirasoro1993,CrisantiSommers1995TAP}.

The expected number of critical points of spherical spin glasses was computed with the Kac--Rice formula by Auffinger, Ben Arous, and \v{C}ern\'y for pure models \cite{AuffingerBenArousCerny2013} and by Auffinger and Ben Arous for mixtures \cite{AuffingerBenArous2013}. For pure models, Subag proved that the number of critical points at low energies concentrates, by a second moment computation \cite{Subag2017complexity}, and Subag and Zeitouni identified the extremal process of critical values \cite{SubagZeitouni2017}. In \cite{Subag2017geometry}, Subag showed that for pure models at sufficiently low temperature, the Gibbs measure concentrates on thin bands around the deepest critical points. Ben Arous, Subag, and Zeitouni obtained a similar description for mixtures close to pure models at low temperature \cite{BenArousSubagZeitouni2020}. Dembo and Subag proved a decomposition of this kind for generic mixtures that are strictly $k$-step replica symmetry breaking at some $\beta>\bs$. For $k=1$, almost all of the Gibbs measure lies in narrow bands around critical points of $H_N$ on the sphere of radius $\sqrt{qN}$, where $q$ is the positive atom of the Parisi measure, and the energies of these critical points are maximal up to $o(N)$ \cite[Proposition~3.1]{DemboSubag2026}. The Dovbysh--Sudakov representation describes the limits of overlap arrays \cite{DovbyshSudakov1982,Panchenko2010DS}. Under hypotheses that ensure the full family of Ghirlanda--Guerra identities \cite{GhirlandaGuerra1998,Panchenko2010GG}, such as generic mixtures or suitable generic perturbations, these limits are ultrametric \cite{Panchenko2013ultrametricity}; see the monograph \cite{Panchenko2013book}. Related stochastic stability identities were derived by Aizenman and Contucci \cite{AizenmanContucci1998}.

\subsection{Dynamics}\label{sec:intro-dynamics}

The rigorous study of Langevin dynamics for spin glasses was initiated by Ben Arous and Guionnet, who proved large deviations for the empirical measure of paths and derived self-consistent limiting equations for soft spins \cite{BenArousGuionnet1995,BenArousGuionnet1997}. Ben Arous, Dembo, and Guionnet proved aging for the spherical $2$-spin model \cite{BenArousDemboGuionnet2001} and derived the Cugliandolo--Kurchan equations for mixed $p$-spin models with a soft spherical constraint and an initial law independent of the disorder \cite{BenArousDemboGuionnet2006}. Dembo, Guionnet, and Mazza treated the hard spherical constraint and analyzed the limiting equations at high temperature \cite{DemboGuionnetMazza2007}. For asymmetric soft-spin dynamics of Sherrington--Kirkpatrick type, Dembo, Lubetzky, and Zeitouni proved universality of the empirical path law at fixed times under independence assumptions on the couplings \cite{DemboLubetzkyZeitouni2021}. Dembo and Subag derived the limiting dynamics for initial conditions spread uniformly over a band of prescribed overlap around a critical point \cite{DemboSubag2020} and for initial conditions drawn from the Gibbs measure, in replica-symmetric models and generic strictly one-step replica symmetry breaking models, in the limit $N\to\infty$ followed by a limit of hard constraint \cite{DemboSubag2025}. For $\bd<\beta<\bs$, where the model is replica symmetric, the limiting correlation of the stationary dynamics has a nonzero plateau at a root of the plateau equation \cite[Proposition~1.14 and Corollary~1.18]{DemboSubag2025}.
These results concern the limit $N\to\infty$ at fixed times, and they do not control the dynamics at times exponential in $N$. Jagannath surveyed the dynamics on short and long time scales \cite{Jagannath2019}.

In finite dimension, Ben Arous, Gheissari, and Jagannath proved differential inequalities for the energy and the gradient along the dynamics, uniformly over initial points \cite{BenArousGheissariJagannath2020}. Sellke studied pure models at sufficiently low temperature. From every initial point, the energy of the dynamics rises above the threshold energy minus any fixed $\eta>0$ and stays above this level up to an exponential time \cite[Theorem~1.1]{Sellke2024}. For an initial law independent of the disorder, the energy stays below the threshold minus a positive $\delta(p,\beta_0)$ on every fixed bounded time interval, uniformly over $\beta\le\beta_0$, with exponentially high probability \cite[Theorem~1.7]{Sellke2024}. Obstructions based on the overlap gap property were proved for approximate message passing \cite{GamarnikJagannath2021} and for Langevin dynamics started independently of the disorder, on the time scales considered there \cite{GamarnikJagannathWein2024}. Huang and Sellke proved sharp hardness results for Lipschitz optimization algorithms, with applications to Langevin dynamics on bounded time intervals \cite{HuangSellke2025}. The energy bounds from independent starts and the hardness results cited here concern bounded time intervals; Sellke's lower bound on the energy from arbitrary starts persists to exponential times.

Exponentially slow relaxation has been studied through the spectral gap. Gheissari and Jagannath proved for pure models with $p\ge3$ that the spectral gap is bounded below at sufficiently high temperature and exponentially small at sufficiently low temperature \cite[Theorem~1]{GheissariJagannath2019}. Ben Arous and Jagannath showed that a free-energy barrier in the overlap implies an exponentially small spectral gap, and verified such a barrier for classes of mixtures that include pure even models with $p\ge4$ \cite[Theorems~1.24--1.26]{BenArousJagannath2018}. In \cite{BenArousJagannath2024}, for pure spherical models with $p\ge4$ at temperatures below a threshold $T_{\mathrm{BBM}}$, they constructed exponentially many bands around critical points such that the dynamics started from the Gibbs measure conditioned on a band has exponentially small probability of escaping before an exponential time. They also formulated a weak notion of shattering and proved it for pure models with $p\ge4$ in an interval of temperatures just above the static transition. Such bands may carry a small part of the Gibbs measure. On the high-temperature side, Kim ruled out shattering in this fixed-overlap critical-point-band sense for all overlaps in the pure cubic model when $\beta\le\bd(3)$, and for specified overlap ranges when $p\ge4$ \cite[Theorem~1.1 and Corollary~1.2]{Kim2026nonshattering}. This concerns the definition of Ben Arous and Jagannath, rather than the stronger cluster decomposition below; it does not establish rapid mixing. El Alaoui, Montanari, and Sellke proved shattering of pure models in a strong sense, with consequences for the spectrum and the equilibrium correlation \cite{ElAlaouiMontanariSellke2025}, which we discuss in \cref{sec:intro-comparison}. Auffinger, El Alaoui, and Sellke studied mixtures with a discontinuous replica symmetry breaking transition and with $\bs<\xi''(0)^{-1/2}$. For all $\beta$ in an interval immediately below $\bs$, they proved that the constrained pressure is strictly increasing on an interval of positive overlaps \cite[Theorem~1.1]{AuffingerElAlaouiSellke2025}, and they deduced an exponentially small spectral gap \cite[Remark~1.2]{AuffingerElAlaouiSellke2025}. Sellke studied spherical models near their ground states. For every mixture, he showed that at all approximate ground states, the gradient is close to an explicit multiple of the point and the bulk of the spectrum of the Hessian has explicit edges, and for even mixtures, that the largest eigenvalue of the Hessian is at most the upper edge, up to a small error \cite[Proposition~1.2 and Theorem~1.3]{Sellke2025marginal}. For even mixtures whose approximate ground states are not marginally stable (in his terms, without full RSB endpoint behavior), he showed by the argument of \cite[Corollary~2.9]{ElAlaouiMontanariSellke2025} that at sufficiently large $\beta$, the stationary dynamics started from a Gibbs sample keeps overlap at least $1-\eps$ with its starting point up to time $e^{cN}$, with probability at least $1-e^{-cN}$, and he noted that this implies exponentially slow mixing \cite[Corollary~1.6]{Sellke2025marginal}. We compare his estimates near the ground state with ours before \cref{prop:mo-all-high}. For finite even mixtures with minimum degree at least four, under a strict separation between the ground-state energy and the energies at which the annealed complexities of critical points of index one and two vanish, Badiei Khuzani proved one-sided lower bounds on the relaxation time, in probability over the disorder: a sequential Arrhenius bound, in which $\beta\to\infty$ at fixed $N$ before $N\to\infty$; a bound at every fixed $\beta$ at which the free energy exceeds $\beta$ times the first of these two energies; and an aggregate Eyring--Kramers bound at every sufficiently large fixed $\beta$ \cite[Definition~1 and Theorems~8--10]{BadieiKhuzani2026}. The current version of that paper notes that these bounds do not show slow relaxation from the uniform law, since the slow test function in the bound at fixed temperature is odd under $x\mapsto-x$, while the uniform law, the Gibbs measure, and the generator are even \cite[Section~4.3]{BadieiKhuzani2026}; compare \cref{sec:intro-starts}. For Ising models, exponentially slow worst-case mixing at low temperature was proved by Sellke for the Sherrington--Kirkpatrick model \cite{Sellke2025} and by Kouraich and Warzel for $p$-spin models \cite{KouraichWarzel2026}.

At high temperature, sampling and fast mixing are now well understood. Huang, Montanari, and Pham gave a polynomial-time algorithm, based on stochastic localization, that samples from the Gibbs measure with vanishing error in total variation when $\beta^2\xi''(q)<(1-q)^{-2}$ for all $q\in[0,1)$ \cite{HuangMontanariPham2024}. Huang, Mohanty, Rajaraman, and Wu proved that simulated annealing with Langevin dynamics samples from the Gibbs measure up to the same stochastic localization threshold \cite{HuangMohantyRajaramanWu2025}. Their algorithm starts from the uniform law and runs through a sequence of temperatures, and weak Poincar\'e inequalities control its warm starts. In a different regime, Huang and Sellke proved mixing in logarithmic time at sufficiently low temperature for strongly topologically trivial landscapes with a unique local maximum, which include suitable external fields \cite{HuangSellke2026}. For Ising spins, spectral gaps and fast mixing at high temperature were proved in \cite{BauerschmidtBodineau2019,EldanKoehlerZeitouni2022,AdhikariBrenneckeXuYau2024,AnariJainKoehlerPhamVuong2024}.

The relation between metastable sets and small eigenvalues of reversible dynamics has a long history. It goes back to the exit and cycle theory of Freidlin and Wentzell \cite{FreidlinWentzell2012} and to Witten's spectral approach to Morse theory \cite{Witten1982} and the semiclassical analysis of multiwell operators \cite{HelfferSjostrand1984}. It was developed in the potential-theoretic approach of Bovier, Eckhoff, Gayrard, and Klein \cite{BovierEckhoffGayrardKlein2000,BovierEckhoffGayrardKlein2001,BovierEckhoffGayrardKlein2002,BovierEckhoffGayrardKlein2004,BovierGayrardKlein2005} and in the work of Helffer, Klein, and Nier \cite{HelfferKleinNier2004}; see the monograph \cite{BovierDenHollander2015}. The works on reversible diffusions and the semiclassical works give sharp asymptotics as a noise parameter tends to zero, while the works on Markov chains and disordered mean-field models also treat growing systems. Our results concern spherical spin glasses in high dimension, with exponentially many slow modes.

\subsection{Comparison with the closest results}\label{sec:intro-comparison}

We begin with the dynamics started from the uniform law. For pure spherical $p$-spin models, Ben Arous and Jagannath conjectured that at every temperature below the dynamical transition, the dynamics started from the uniform law takes exponential time to thermalize, with probability tending to one, and that it thermalizes in bounded time above the dynamical transition \cite[arXiv v1, Section~2.5, Conjecture~6]{BenArousJagannath2024}.
Their thermalization time is the first time $t$ at which $\max_{\norm{f}_\infty\le1}\int(P_tf-\pi(f))^2\dd\nu_N\le1/e$. Taking $f=2\ind_A-1$ and applying the Cauchy--Schwarz inequality, we see that this maximum is at least $4\norm{\nu_NP_t-\pi}_{\TV}^2$. Under the hypotheses of \cref{thm:main-condensed}, this thermalization time is therefore at least $e^{\gamma N}$ with probability tending to one, in the normalization of \cite{BenArousJagannath2024} in which the dynamics is reversible with respect to the Gibbs measure, after the corresponding constant rescaling of time. (The generator displayed in \cite[p.~2 of arXiv v1]{BenArousJagannath2024} and the stated Gibbs measure differ by a factor of two in the drift; we compare temperatures through the inverse temperature of the Gibbs measure.) In particular, for pure models with $p\ge4$, the lower bound of the conjecture holds at every temperature below the static transition. For the mixtures that satisfy the 1RSB assumptions of \cref{def:intro-certificate}, the same lower bound follows, although the conjecture was stated for pure models. \Cref{thm:main-condensed} does not resolve the conjecture: the conjecture also concerns the interval $\bd<\beta\le\bs$ and the case $p=3$, where our methods do not apply to the uniform start, and it predicts thermalization in bounded time above the dynamical transition. On bounded time intervals, and at temperatures low enough that the Gibbs measure concentrates above the threshold energy, the energy bound of Sellke \cite[Theorem~1.7]{Sellke2024} separates the dynamics of pure models, started independently of the disorder, from equilibrium. To our knowledge, \cref{thm:main-condensed} is the first result that separates the uniform start from equilibrium in total variation up to times exponential in $N$ throughout the condensed phase, for pure spherical models with $p\ge4$ at every $\beta>\bs$, and for the mixtures that satisfy the 1RSB assumptions. The spectral gap estimates above concern worst-case starting points, and, as the example of \cite{GheissariGrantYu2026} shows, they do not imply such a statement.

The closest antecedent of \cref{thm:main-memory} is the work of El Alaoui, Montanari, and Sellke \cite{ElAlaouiMontanariSellke2025}. For pure models, they proved that with probability at least $1-e^{-cN}$, all but an exponentially small part of the Gibbs measure is covered by exponentially many clusters of small diameter and exponentially small mass, separated by bottlenecks. This holds for sufficiently large $p$ and every $\beta$ in an interval $(C,\bs(p))$, where $C$ is an absolute constant \cite[Theorem~2.1]{ElAlaouiMontanariSellke2025}, and for $p\ge3$ whenever the constrained pressure is strictly increasing on an interval of overlaps above $0.9999$ \cite[Theorem~2.4]{ElAlaouiMontanariSellke2025}. They deduced that the first $e^{cN}$ eigenvalues of $-L$ are at most $e^{-cN}$ \cite[Corollary~2.8]{ElAlaouiMontanariSellke2025}, and that, with probability at least $1-e^{-cN}$ jointly over the disorder, the stationary initial point, and the dynamics, $\inf_{0\le t\le e^{cN}}R(X_0,X_t)\ge1-\Delta^2$, where $\Delta$ is defined from the scales of the clusters \cite[Corollary~2.9]{ElAlaouiMontanariSellke2025}; in particular, their annealed correlation is at least $1-2\Delta^2$ on this interval.
Exponentially many slow modes and an exponential-time plateau were therefore known for pure models in these regimes. By the same argument, Sellke obtained such a plateau at sufficiently large $\beta$ for even mixtures whose approximate ground states are not marginally stable \cite[Corollary~1.6]{Sellke2025marginal}. \Cref{thm:main-memory} extends them in three ways. First, it applies to every finite mixture, including a quadratic term and odd degrees, and to every $\beta\in(\bd,\bs)$. Second, the plateau level is given in terms of $q_*$: every level below $q_*^2$ with probability at least $1-Ce^{-aN}$, and every level below $q_*$ with probability tending to one (\cref{thm:main-plateau}). Third, it adds the conclusions on Gibbs-typical starting points and on stationary information. For finite and infinite mixtures, \cref{thm:main-cusp} gives eigenvalue bounds close to $\bd$. On the other hand, where the results of \cite{ElAlaouiMontanariSellke2025} apply, they contain more: an explicit decomposition into clusters and a pathwise plateau. In their regime, trapping from Gibbs-typical points and information at least linear in $N$ can also be derived from the persistence of clusters. We construct no clusters, and our argument is different (see \cref{sec:intro-ideas}).

In the limit $N\to\infty$ at fixed times, followed by a limit of hard constraint, the plateau of the stationary correlation for $\bd<\beta<\bs$, where the model is replica symmetric, was identified by Dembo and Subag \cite[Proposition~1.14 and Corollary~1.18]{DemboSubag2025}. \Cref{thm:main-memory,thm:main-plateau} hold in finite dimension up to exponential times, but they are only lower bounds. In their regimes, \cite{BenArousJagannath2018,AuffingerElAlaouiSellke2025} show that one eigenvalue is exponentially small; in \cref{thm:main-memory}, we obtain exponentially many for every $\beta\in(\bd,\bs)$, together with the plateau. The bands of \cite{BenArousJagannath2024} trap the dynamics started from the Gibbs measure conditioned on a band; since they may carry a small part of the Gibbs measure, they do not lead to a plateau of the equilibrium correlation, as noted in \cite{ElAlaouiMontanariSellke2025}.

\section{Ideas of the proofs}\label{sec:intro-ideas}

A detailed overview of the proofs is given in \cref{sec:overview}. Here we describe the main steps and indicate where the tools come from.

\paragraph{The condensed phase.}
Let $(\xi,\beta)$ be strictly 1RSB with parameters $x,Q$ and a nondegenerate contact at $Q$. The proof of \cref{thm:main-condensed} rests on a description of the Gibbs measure that goes back to the work of Subag on pure models at low temperature \cite{Subag2017geometry}. Consider the center field $H_Q(u)=H_N(\sqrt Q\,u)$ on $S_N$. Its critical points with value in a window of fixed width around a deterministic centering $m_N=NE_0+O(\log N)$, and with normalized radial derivative $K_Q(u)/N$ close to an explicit value $K_0$, have bounded expected number, and the maximum of $H_Q$ lies in this window with probability close to one when the width is large. The bands $\{x\in S_N:|R(x,u)-\sqrt Q|\le b\}$ around these critical points $u$ carry almost all of the Gibbs measure. On the other hand, a point drawn from $\nu_N$ has overlap of order $N^{-1/2}$ with each of these points. We show that before time $e^{\gamma N}$, the diffusion does not reach the caps $\{x:|R(x,u)|\ge a\}$, which contain the bands when $a<\sqrt Q-b$. The union of the bands is then a set of Gibbs measure close to one that the diffusion does not visit. In \cref{sec:reduction}, we turn these three properties into \cref{thm:main-condensed}, using the contraction of total variation and a first moment computation with the Kac--Rice formula. We choose the exponent $\gamma$ before the width of the window. The inputs are proved as follows.

\emph{Statics} (\cref{sec:thermal}). We identify the Parisi measure under the 1RSB assumptions through the Crisanti--Sommers optimality criterion \cite[Proposition~2.1]{Talagrand2006}, \cite[Lemma~1.2]{HuangSellke2024}, and the stationarity equations for two atoms are the classical one-step equations of Crisanti and Sommers \cite{CrisantiSommers1992}. When $\xi''^{-1/2}$ is strictly convex, we verify the 1RSB assumptions at every $\beta>\bs$ by the argument of Crisanti and Sommers that bounds the number of contacts, as presented by Talagrand \cite[Proposition~2.2]{Talagrand2006}; its version at zero temperature is \cite[Lemma~3.21]{HuangSellke2024}. Using the generalized TAP representation of Subag \cite{Subag2024}, we show for every finite mixture and every $\beta$ that an overlap is multisamplable if and only if it belongs to the contact set, that is, it is a zero of the first-variation function of the Parisi measure (\cref{thm:th-contact}). Subag proved that every zero is multisamplable, and the converse for generic mixtures \cite[Theorem~10]{Subag2024}. Under the 1RSB assumptions, the multisamplable overlaps are $0$ and $Q$. The ground state of $H_Q$ is then strictly 1RSB at zero temperature in the sense of \cite[Definition~4]{HuangSellke2024}, by the restriction and rescaling of minimizers at a contact point \cite[Proposition~11]{Subag2024} and the formula for the ground state energy \cite{ChenSen2017,JagannathTobasco2017low}. This gives the ground-state energy $E_0$ of $H_Q$, which determines the centering $m_N$, and the value $K_0$ of the normalized radial derivative.

\emph{The directing measure} (\cref{sec:directing}). Every limit of the overlap array of independent Gibbs samples is generated by a random measure on the unit ball of a Hilbert space, the directing measure \cite{DovbyshSudakov1982,Panchenko2010DS}. We show that it is carried by vectors of squared norm $Q$ (\cref{thm:dir-norm}). We exclude other positive squared norms with the characterization of multisamplable overlaps, and an atom at the origin with a single identity of Ghirlanda--Guerra type, which we derive from the self-averaging argument of Panchenko \cite{Panchenko2010GG}. This argument uses no ultrametricity and no Poisson--Dirichlet structure.

\emph{Counts} (\cref{sec:kacrice,sec:purecover,sec:mixone,sec:mixpair}). We use the Kac--Rice formula \cite{AdlerTaylor2007,AzaisWschebor2009,Stecconi2022,MathisStecconi2024} with weights that are bounded Borel functions of the whole field, such as hitting probabilities of the diffusion. The other inputs are the law of the field at a point of a smaller sphere \cite{BenArousSubagZeitouni2020} and bounds on all moments of the Hessian determinant. For pure models, the bound on the counts follows from a first moment computation and the tightness of the centered maximum \cite{SubagZeitouni2017}. For mixtures, we follow the conditional trial argument of Huang and Sellke \cite[Section~3.3]{HuangSellke2024} to show that a critical point near the ground state has no other high-energy points at intermediate overlap. We then bound the second moment of the number of critical points in a fixed window, in the spirit of \cite{Subag2017complexity}. The proof uses a local factorization of the two-point Kac--Rice integrand and an elementary bound for spherical codes in a cap \cite{Rankin1955,BargMusin2007}.

\emph{Concentration on bands} (\cref{sec:purecover,sec:mixbands}). Using the estimates of Subag for replicated bands with a fixed number of replicas \cite[Lemmas~17--19 and Proposition~24]{Subag2024}, together with the directing measure, we place a Gibbs sample in a band around a critical point of nearly maximal value. We then pass to critical points in a fixed window, with Laplace estimates for the partition functions of bands and the lower tail of normalized partition functions \cite[Lemma~3.2 and Remark~3.3]{HuangMontanariPham2024}. For pure models at sufficiently low temperature, the concentration of the Gibbs measure on bands around the deepest critical points was proved by Subag \cite[Theorem~1]{Subag2017geometry}, and for generic mixtures that are strictly one-step replica symmetry breaking, the decomposition of Dembo and Subag places almost all of the Gibbs mass in narrow bands around critical points (see \cref{sec:intro-statics}). We prove that the bands around the critical points in a fixed window carry almost all of the Gibbs mass, for pure models with $p\ge4$ at every $\beta>\bs$ and for all mixtures that satisfy the 1RSB assumptions. Compared with \cite[Proposition~3.1]{DemboSubag2026}, our mixtures are finite, whereas generic mixtures have infinitely many nonzero coefficients, and the centers lie in an energy window of fixed width, which contains a bounded expected number of critical points, instead of a window of width $o(N)$. On the other hand, their bands are narrower, and their decomposition also describes the Gibbs weights of the bands.

\emph{The hitting estimate} (\cref{sec:entrance}). Fix a point $u$ and condition the field to have a critical point of $H_Q$ at $u$ with prescribed value and radial derivative. Using Gaussian regression, we write the conditioned field as $H_N$ plus a correction that is of order $N$ near $u$ but vanishes to order $|R(x,u)|^d$ near the equator of $u$. The diffusion is the original Langevin diffusion in the conditioned field. We compare it, for an axis $u$ that is uniform on $S_N$, with the diffusion in the unconditioned field, for which the Gibbs mass of the caps around $u$ is exponentially small by rotation invariance. For the averages over latitudes in this comparison, we use the logarithmic Sobolev inequality on the sphere, through the curvature criterion of Bakry and \'Emery \cite{BakryEmery1985}, and the Funk--Hecke formula \cite{AtkinsonHan2012}. The main obstacle is the length $e^{\gamma N}$ of the time interval, since a comparison of path laws costs an amount proportional to the time. To overcome it, we compare resolvents instead, at a cost controlled by the relative entropy of $\nu_N$ with respect to the Gibbs measure, which is of order $N$, and not by the time (\cref{lem:ent-resolvent}). The resulting error $N^{1-d/2}+N^{3-d}$ tends to zero because $d\ge4$, and this is where the hypothesis on the minimum degree enters.

\emph{Determinant weights and other initial laws} (\cref{sec:kacrice,sec:reduction,sec:condproof}). The hitting estimate for the caps under critical-point conditioning concerns ordinary Gaussian conditioning, while the Kac--Rice formula weights each critical point by the absolute determinant of its Hessian. We transfer the estimate with H\"older's inequality and the moment bounds for the determinant, as in the conditional Cauchy--Schwarz argument of \cite[Corollary~3.16]{HuangSellke2024}. The transfer turns the bound $N^{-(d-2)/2}$ into $N^{-(d-2)(1-1/r)/2}$ for a fixed $r>1$, and this is why the endpoint $\kappa=(d-2)/2$ is excluded in \cref{thm:main-initial}(ii). Part (i) of \cref{thm:main-initial} is an instance of the classical symmetrization over a compact group, the same averaging that underlies the Hunt--Stein theorem in statistics \cite[Section~8.5]{LehmannRomano2005}: the expected distance to equilibrium is convex and rotation invariant as a function of the initial law, and $\nu_N$ is the rotation average of every initial law, so $\nu_N$ minimizes the expected distance. For parts (ii) and (iii), we combine a polynomial bound on the mass of the bands at time $e^{\gamma N}$ with a union bound and with data processing for relative entropy \cite[Theorem~2.17]{PolyanskiyWu2025}.

\paragraph{Below the static threshold.}
The proof of \cref{thm:main-memory} has a probabilistic part and a deterministic part. In the probabilistic part, we study the constrained pressure $F_{\mathrm{FP}}(r)$, the free energy per site of the configurations whose overlap with a reference point drawn from $\pi$ is close to $r$, which is the potential of Franz and Parisi \cite{FranzParisi1995} written as a pressure. We compute it through a planted model, in which the reference point is drawn first and the disorder is tilted toward it, following El Alaoui, Montanari, and Sellke, who treated pure models \cite[Propositions~3.7 and~3.8]{ElAlaouiMontanariSellke2025}; planted models go back to \cite{AchlioptasCojaOghlan2008,KrzakalaZdeborova2009}. We transfer exponentially rare events from the planted law to the joint law of the disorder and a reference point drawn from the Gibbs measure, as in \cite[Lemma~3.5]{ElAlaouiMontanariSellke2025} and \cite[Lemma~2.3]{AuffingerElAlaouiSellke2025}. Dembo and Subag also obtained, through Gaussian conditioning, a formula for the constrained pressure of general spherical models in the replica-symmetric phase $\beta\le\bs$, and its extension to generic strictly one-step replica symmetry breaking models in the condensed phase $\beta>\bs$ \cite[Corollary~2.1 and Theorem~2.3]{DemboSubag2026}.
We bound $F_{\mathrm{FP}}$ by an explicit trial pressure, as in \cite[Proposition~3.9]{ElAlaouiMontanariSellke2025}, and show that equality holds at $q_*$ for every finite mixture, using the monotonicity of the ratio $r\xi''(r)/\xi'(r)$. Seen from a Gibbs-typical reference point $x$, the Gibbs measure of a collar $\{u:q_1\le R(x,u)\le q_2\}$ below $q_*$ is then exponentially smaller than that of the cap $\{u:R(x,u)>q_2\}$. Further, every cap of fixed latitude has exponentially small Gibbs mass, by the replica concentration of Huang, Montanari, and Pham \cite[Lemma~7.7]{HuangMontanariPham2024}.

In the deterministic part, we use one observable, the overlap-weighted Gibbs mean $B_N$ (\cref{def:sp-barycenter}). For $x\in S_N$, $B_N(x)$ is the mean of $U$ under the Gibbs measure reweighted by a cutoff that suppresses the points $u$ with $R(x,u)<q_1$. Its norm is at most one, its Dirichlet energy is exponentially small because the collar has exponentially smaller mass than the cap, and $U\cdot B_N\ge q_1-o(1)$ on most of the Gibbs measure. We count slow modes with the averaged variational principle of Harrell and Stubbe \cite[Theorem~3.1]{HarrellStubbe2014} and of El Soufi, Harrell, Ilias, and Stubbe \cite[Theorem~2.1]{ElSoufiHarrellIliasStubbe2017}, for which Kr\"oger's averaged trial functions are an antecedent \cite{Kroger1992}. We apply it to translated bumps of $B_N$, whose overlaps are controlled by the Gibbs mass of caps. In \cite{ElAlaouiMontanariSellke2025}, the eigenvalues are counted with smoothed indicators of clusters with disjoint supports; the averaged principle needs no disjointness. By spectral calculus, we then obtain the projection and the correlation bounds (ii) and (iii) at levels below $q_1^2$, and $q_1$ can be taken close to $q_*$. We deduce (iv) from a bound on the mean-square displacement of $B_N$ along the dynamics, Markov's inequality, and binary data processing, and (v) from the chain identity for mutual information \cite[Theorem~4.1]{PolyanskiyWu2025}.

For \cref{thm:main-plateau}, we note that $|B_N(x)|^2$ is the mean overlap of two independent samples from the reweighted Gibbs measure. We show that this overlap concentrates at $q_*$ by differentiating a transverse pressure with respect to the amplitudes of its covariance, a classical device in the theory of mean-field spin glasses \cite{Panchenko2008}, \cite[Proposition~2.2]{AuffingerElAlaouiSellke2025}, combined with the envelope theorem \cite{MilgromSegal2002} and the convergence of derivatives of convex functions \cite[Theorem~24.5]{Rockafellar1970}. The transfer from the planted law uses contiguity \cite[Lemma~3.2]{HuangMontanariPham2024}, and this step has no rate. For \cref{thm:main-cusp}, we expand $A$ near $q_d$ to compute the barrier of the trial pressure. We also use a bound on the Gibbs mass of caps whose exponent is uniform in the temperature. For infinite mixtures, we use bounds on the derivatives of homogeneous polynomials that follow from Bernstein's inequality \cite{QueffelecZarouf2019}.

The tools above are classical or published. Our proofs establish or combine the following ingredients: the hitting estimate for the caps under critical-point conditioning and its comparison of resolvents; the concentration of the Gibbs measure on bands around critical points in a fixed window, for pure models with $p\ge4$ at every $\beta>\bs$ and for mixtures that satisfy the 1RSB assumptions, together with the fact that the directing measure is carried by vectors of squared norm $Q$; the identification of the multisamplable overlaps with the zeros of the first-variation function for every finite mixture; the equality between the constrained pressure and the trial pressure at $q_*$ for general mixtures; the overlap-weighted Gibbs mean and its use with the averaged variational principle; and the proof of \cref{thm:main-plateau} through derivatives in the amplitudes.

\section{Organization of the paper}\label{sec:intro-organization}

\Cref{sec:overview} gives an overview of the proofs. \Cref{part:prelim} consists of \cref{sec:prelim}, which fixes conventions for constants, for statements that hold with high probability, and for the order in which parameters are chosen (\cref{sec:pre-conventions}), and collects facts on the model, the dynamics, Markov semigroups, and statics that both halves of the paper use. The two halves are independent of each other; a reader interested in one of them can skip the other.

\Cref{part:condensed} proves \cref{thm:main-condensed,cor:main-classes,thm:main-initial}. \Cref{sec:reduction} reduces separation to three properties of critical points and bands, and treats other initial laws. \Cref{sec:thermal} derives the static consequences of the 1RSB assumptions and verifies these assumptions for the classes of \cref{cor:main-classes}, with finite computations for example (d) in \cref{app:certificate}. \Cref{sec:directing} proves that the directing measure is carried by vectors of squared norm $Q$. \Cref{sec:kacrice} contains the Kac--Rice formula and moment bounds for determinants, and \cref{sec:entrance} proves the conditioned hitting estimate. \Cref{sec:purecover} proves bounds on counts and the concentration of the Gibbs measure on bands for pure models, and \cref{sec:mixone,sec:mixpair,sec:mixbands} do so for mixtures. \Cref{sec:condproof} assembles the proofs.

\Cref{part:prestatic} proves \cref{thm:main-memory,thm:main-plateau,thm:main-cusp}. \Cref{sec:pressure} studies the constrained pressure, \cref{sec:spectral} contains deterministic estimates for counting slow modes and for the dynamics started from Gibbs-typical points, together with the overlap-weighted Gibbs mean, and \cref{sec:memory} proves \cref{thm:main-memory}. \Cref{sec:plateau} proves \cref{thm:main-plateau}, \cref{sec:infinite} develops the infinite-mixture inputs and proves \cref{thm:main-cusp}. Notation specific to a section is introduced at its start.
